\documentclass[a5paper,11pt]{book}

% {{{ latex things

\title{\textbf{EX VINUM}}
\author{by tsbohc}
\date{}

\usepackage{longtable}
\usepackage{multicol}
\usepackage{CormorantGaramond}
\usepackage{tabularx} % page width
\usepackage{calc}
\usepackage{booktabs}
\usepackage[margin=0.6in, bottom=0.7in]{geometry}
\usepackage{array}
\usepackage[table]{xcolor} % table row colors
\usepackage{hyperref} % links
\usepackage{blindtext} % lorem ipsum
\usepackage[utf8]{inputenc}
\usepackage[russian]{babel}
\usepackage{titlesec} % titles customization
\usepackage{epigraph}
%\usepackage{graphicx}
%\usepackage{background}
\usepackage{lettrine} % big letters in the beginning of the chapter
%\usepackage{indentfirst}
\usepackage{fancyhdr} % page numbers
\usepackage{tocbasic}

%\hyphenpenalty=100000

% messing with layout
%\raggedbottom
%\widowpenalty=1000
%\clubpenalty=1000

%\backgroundsetup{scale=0.75,angle=0,contents=\includegraphics{paper.jpg}}

% set page numbers at the bottom
\pagestyle{fancy}
\fancyhf{}
\renewcommand{\headrulewidth}{0pt}
\cfoot{\thepage}

% change space between columns
\setlength{\columnsep}{18pt}

% remove epigraph line 
\setlength\epigraphrule{0pt}

% just table things
\definecolor{lightgray}{gray}{0.95}
\definecolor{white}{gray}{1}
\newcolumntype{Y}{>{\centering\arraybackslash}X}
\newcolumntype{Z}{>{\ \ \ }l}
\newcolumntype{P}[1]{>{\centering\arraybackslash}p{#1}}
\setlength{\extrarowheight}{2pt}

% toc chapters
\setcounter{tocdepth}{1}

\addto\captionsrussian{% Replace "english" with the language you use
  \renewcommand{\contentsname}%
    {Table Of Contents}%
}

\DeclareTOCStyleEntry[
linefill=\bfseries\TOCLineLeaderFill,beforeskip=2pt
]{tocline}{chapter}
\newcommand\chapterprefixintoc[1]{\MakeUppercase{\chaptername}~#1~}

% titles
\titlespacing*{\chapter}{0pt}{12pt plus 4pt minus 4pt}{8pt plus 1pt minus 1pt}
\titlespacing*{\section}{0pt}{12pt plus 4pt minus 4pt}{10pt plus 1pt minus 1pt}
\titlespacing*{\subsection}{0pt}{8pt plus 2pt minus 2pt}{4pt plus 0pt minus 0pt}

\titleclass{\chapter}{straight}

%\definecolor{reddish}{RGB}{76,28,12}
\definecolor{reddish}{RGB}{0,0,0}

\titleformat{\chapter}
  {\color{reddish}\huge\bfseries\scshape} % format
  {}                % label
  {0pt}             % sep
  {} % before-code
  [{\titlerule[0.1pt]}]

\titleformat{\section}
  {\color{reddish}\large\bfseries\scshape} % format
  {}                % label
  {0pt}             % sep
  {\Large}          % before-code
  [{\titlerule[0.1pt]}]

\titleformat{\subsection}
  {\color{reddish}\large\bfseries\scshape} % format
  {}                % label
  {0pt}             % sep
  {\large}          % before-code

\setlength{\DefaultNindent}{0pt}
\setlength{\DefaultFindent}{6pt}

\begin{document}

\maketitle
\tableofcontents
\pagebreak

% }}}

% {{{ intro
\chapter{Introduction}
\begin{multicols}{2}
\epigraph{\emph{Да раскинутся границы твоего аутизма до дальних цветочных полей и дымкой затянутых горных вершин.}}{--- бухой волшебник}

\lettrine{E}{x vinum} -- настольная ролевая игра, в сердце которой лежат любовь к рпг и слабоумие, границы между которыми были размыты вином. Так она и получила свое название.

Игра делает акцент на ролевом аспекте, упрощая, но при этом не отказываясь от признанных систем.

Мир игры -- классический фэнтези в стиле D\&D, не претендующий на строгий реализм.

\section{Gameplay}
Dungeon Master описывает ситуацию. Детали окружения, персонажей. Однако, далеко не вся информация доступна игрокам сразу.

Игроки описывают, что они хотят сделать. Во время боя это происходит по порядку очередности ходов. Вне боя, игроки могут действовать свободно -- один может взламывать сундук, а другой стоять на стороже.

Иногда решить проблему легко, например открыть дверь. В таком случае DM может просто сказать, что дверь открылась, и описать, что лежит за ней. В других случаях, успешность действий игроков определяется бросками костей.

DM описывает результаты действий игроков, и на этом цикл повторяется.

Стоит помнить, что это не компьютерная игра. Действия участников ограничиваются лишь их воображением и фэнтезийной версией законов физики, а минмаксинг может обернуться плачевно.

\section{Nota Bene}
DM и игроки работают вместе, чтобы создать общий нарратив, а не играют друг против друга.

Что бы ни было написано в правилах, DM всегда имеет последнее слово.

DM должен поощрять изобретательность игроков их способность вживаться в роли, даже если их идеи не реалистичны.

Главная цель не победить, а хорошо провести время в кругу друзей.

% }}}

% {{{ races 
\end{multicols}
\chapter{Races}
\begin{multicols}{2}

\lettrine{Н}{}иже перечислены разумные, игровые расы ex vinum. Полный список существ далеко не закачивается этим разделом.

\section{Human}
Ну, вы поняли. Люди преуспевают во многом, но реже добиваются невероятных результатов. Являются самой многочисленной расой, включающей в себя огромное количество разных культур.

\section{Dwarf}

Воспитанные грубой средой обитания и тяжелой работой, Дворфы ценят свое мастерство. Чаще всего они становятся ремеслeнниками или шахтерами. Любят провести вечер пересчитывая свое золото за парой-тройкой-десятком кружек эля.

Невысокого роста. Тучные, мускулистые, с бурым или темным цветом волос, грубоватым голосом и громким смехом.

%Сильные и выносливые, но слегка неповоротливые.

\section{Elf}

Изящные и благоразумные, Эльфы тесно связаны с магией этого мира. Часто становятся писателями или художниками. Большую часть времени проводят наедине с природой и книгами.

Ростом чуть ниже человека. Стройные, светловолосые, часто с голубым или зеленым оттенком глаз. Голоса мелодичны и благозвучны.

%Умные и харизматичные, но не обладают большой физической силой.

\section{Tiefling}

Тифлинги несут за собой бремя своей родословной -- они потомки людей и дьявола. Их полу-демоническое обличие не пробуждают в людях теплых чувств, лишь только косые взгляды. Тифлинги, оказавшиеся в этом мире не по своей воли, делают все, чтобы выжить.

Ростом с человека. Цвета кожи те же, что и у людей, но могут принимать и красноватые оттенки. Длииинный хвост. На голове большие рога. Туманные глаза, темные волосы.

%\smallskip
%+2 dexterity, -1 intelligence

%\section{Duergar}
%
%\smallskip
%+1 strength, +1 dexterity, -1 intelligence

\section{Drow}

Темные Эльфы произошли от дверней подрасы Эльфов, изгнанных обитать в глубинах земли. Те, что выбираются на поверхность, нередко идут по пути зла, но встречаются и исключения. По большей части становятся путешественниками или наемниками.

Чуть меньше и стройнее своих собратьев. Обладают темным цветом кожи, практически белым цветом волос и бледными глазами с оттенком.

%\smallskip
%-1 strength, +1 dexterity, +1 intelligence

\section{Aasimar}

Азимары подобны Тифлингам в своем происхождении, но в их жилах течет божественная кровь.  Азимары от природы расположены к благотелям, а в своих снах ведут разговоры со своим божеством. Однако, предпочитают не раскрывать свою родословную.

Выше человека. В цветах волосы и кожи встречаются как человеческие, так и более металлические оттенки. Другими рассамы описываются как люди невероятной красоты. 


%\smallskip
%+1 strength, -1 dexterity, +1 intelligence

% }}}

% {{{ roleplaying

\end{multicols}
\chapter{Roleplaying}
\begin{multicols}{2}
\lettrine{К}{}ем будешь ты? Добродушным дворфом с раскатистым смехом, знающим лучшие таверны в городе? Или тифлингом, чьи руки в крови, потому что единственным выбором была смерть?

\section{Race}
Выбор расы влияет не только на физический облик и характеристики, но и ваше положение в мире, а также отношения с другими рассами. Кроме этого, ваша история вытекает из происхождения.

\section{Age}
Некоторые расы, например дворфы и эльфы, живут на несколько сотен лет дольше чем люди. В результате чего имеют другую перспективу на мир.

\section{Sex}
Пол обладает большим влиянием на ваш характер и прошлое. Не все женщины покинув родной дом тут же становятся рыцарями. Рыцарями становятся те женщины, чьим сердцам дорого то, что они хотят защитить.

\section{Class}
Каждый искатель приключений обладает той или иной профессией. Профессия определяет, владаеете ли вы магией или хорошо управляетесь с оружием, и какую роль вам предстоит выполнять в кампании.

\section{Personality}
Личность составляют идеи которые вами движат и недостатки которыми обладаете. Ваши сильные и слабые стороны, ваш характер, и то, как вы ведете себя с другими людьми.

\section{Backstory}
Какой была ваша жизнь до начала кампании? Может быть, вас отправили на важное задание, или вовсе изгнали? Как вы преобрели свои навыки и стали своим классом? Что вы хотите защитить или чем желаете овладеть? Есть ли у вас возлюбленный или возлюбленная? Потеряли ли вы старого друга или преследуете врага?

\section{Appearance}
Все, что связано с внешним обликом. Телесложение, цвет кожи, волос, глаз. Шрамы, полученные в бою и татуировки. Детали одежды и украшения, талисманы.

% }}}

% {{{ character creation

\end{multicols}
\chapter{Character Creation}
\begin{multicols}{2}

\lettrine{С}{}ейчас хорошее время, чтобы отложить эту книгу и подумать за кого вы хотите играть. Встретившись за столом, вам предстоит рассказать о своем персонаже и познакомиться с другими участниками группы.

Один совет -- попробуйте себя в роли, не совсем похожей на вас -- но такой, которую вам интересно было бы исполнять.

\section{Dice}

Броски костей обозначаются двумя числами через букву d. Например, 3d6 означает три броска кости с 6 сторонами. В случае отсутствия других указаний, подразумевается сумма указанных бросков.

В случае с 2d8+3 к сумме двух бросков d8 прибавляется 3.

В основе ролевой системы ex vinum лежит d12. Это обусловлено желанием добавить динамики традиционной системе d20.

\section{Attributes}
Атрибуты выражают физические и умственные способности персонажа. Определяются численными значениями, в среднем от 3 до 12.

Например, персонаж c 3 интеллекта не способен читать, крайне рассеян и может заговариться.

\subsection{Strength}
Определяет физическую силу и выносливость. Имеет большое влияние на общий запас здоровья.

\subsection{Dexterity}
Определяет ловкость и скорость реакции, а также внимательность. Слабо влияет на запас здоровья.

\subsection{Intelligence}
Определяет ум и сообразительность, фоновые знания, а также уменее вести разговор. Не влияет на запас здоровья.

\section{Rolling Attributes}
Чтобы определить свои базовые атрибуты, игрок бросает 3d4 и складывает полученные значения. Затем повторяет этот процесс еще два раза. Полученные три ролла игрок распределяет на свое усмотрение.

Перед началом игры DM может предложить свой метод определения атрибутов.

\subsection{Best Backstory Award}
Самому интересному персонажу по мнению большинства DM предоставит +1 к выбранному им атрибуту.

\subsection{Racial Bonuses}
На данном этапе к базовым атрибутам применяются расовые бонусы.

\smallskip
\noindent
\rowcolors{1}{}{lightgray}
\setlength\tabcolsep{0pt}
\begin{tabularx}{\linewidth}{@{}ZYYY@{}}
  \textbf{race} & \textbf{str} & \textbf{dex} & \textbf{int} \\
  \hline
  dwarf     & +2 & -1 &    \\
  elf       & -1 &    & +2 \\
  tiefling  &    & +2 & -1 \\
  %duergar   & +1 & +1 & -1 \\
  drow      & -1 & +1 & +1 \\
  aasimar   & +1 & -1 & +1
\end{tabularx}
\smallskip

Люди прибавляют +1 к любому выбранному атрибуту.

\subsection{Health}

Максимальный запас здоровья расчитывается по следующей формуле: \textbf{6 × strength + 2 × dexterity}.

\section{Attribute Modifiers}
Моды используются чтобы подчеркнуть, насколько хорошо или плохо персонаж владеет тем или иным навыком, и участвуют в подсчете роллов.

Высокие значения атрибутов пораждают положительные моды, низкие -- отрицательные.

Расчитать свои моды можно обратившись к таблице:

\smallskip
\noindent
\rowcolors{1}{}{lightgray}
\setlength\tabcolsep{0pt}
\begin{tabularx}{\linewidth}{@{}YYYY@{}}
  \textbf{attr} & \textbf{mod} & \textbf{attr} & \textbf{mod} \\
  \hline
  1-2 & -3 & 9-10 & +1 \\
  3-4 & -2 & 11-12 & +2 \\
  5-6 & -1 & 13-14 & +3 \\
  7-8 & -0 & 15-16 & +4
\end{tabularx}

% }}}

% {{{ character sheet 

\end{multicols}
\chapter{Character Example}

\setlength{\arrayrulewidth}{0.1pt}

\noindent
\setlength\tabcolsep{0pt}
\def\arraystretch{1.3}
\begin{tabularx}{\linewidth}{@{}YYYYYYYYYYYYY@{}}

  \multicolumn{13}{c}{\large\sc\textbf{Name:}} \\ 
  \multicolumn{13}{c}{\Large\emph{Ailon Dawnguard}} \\ \cline{4-10}

  \addlinespace[0.5cm]

  \multicolumn{6}{c}{\large\sc\textbf{Attributes:}} &
  \multicolumn{1}{c}{} &
  \multicolumn{6}{c}{\large\sc\textbf{Character:}}
  \\ \cline{1-6}

  \multicolumn{2}{c|}{\small\sc\textbf{Str}} &
  \multicolumn{2}{c|}{\small\sc\textbf{Dex}} &
  \multicolumn{2}{c}{\small\sc\textbf{Int}} &
  \multicolumn{1}{c}{} &
  \multicolumn{6}{Z}{\emph{Полу-эльф Волшебник, 24, ж.}}
  \\ \cline{8-13}

  \multicolumn{2}{c|}{\LARGE{+0}} &
  \multicolumn{2}{c|}{\LARGE{+1}} &
  \multicolumn{2}{c}{\LARGE{+2}} &
  \multicolumn{1}{c}{} &
  \multicolumn{6}{Z}{\emph{Сбежала из монастыря чтобы}}
  \\ \cline{1-6} \cline{8-13}

  \multicolumn{2}{c|}{\large{7}} & 
  \multicolumn{2}{c|}{\large{9}} &
  \multicolumn{2}{c}{\large{11}} &
  \multicolumn{1}{c}{} &
  \multicolumn{6}{Z}{\emph{увидеть мир таким, какой он}}
  \\ \cline{8-13}

  \multicolumn{6}{c}{\large\sc\textbf{Stats:}}  &
  \multicolumn{1}{c}{} &
  \multicolumn{6}{Z}{\emph{есть. Хочет стать лучше, но}}
  \\ \cline{1-6} \cline{8-13}

  \multicolumn{3}{c|}{\small\sc\textbf{Max Hp}} &
  \multicolumn{3}{c}{\small\sc\textbf{Gold}} &
  \multicolumn{1}{c}{} &
  \multicolumn{6}{Z}{\emph{боится потерять себя. Веснушки,}}
  \\ \cline{8-13}

  \multicolumn{3}{c|}{\LARGE{46}} &
  \multicolumn{3}{c}{\LARGE{75}} &
  \multicolumn{1}{c}{} &
  \multicolumn{6}{Z}{\emph{Русые волосы, невысокий рост.}}
  \\ \cline{1-6} \cline{8-13}
  &  &  &  &  &  &  &  &  &  &  &  & \\
\end{tabularx}

\begin{multicols}{2}

Выше живой пример заполненного листа персонажа, на примере моего протагониста \emph{Baldur's Gate 2}.

Игрокам будут выданы дополнительные листки для заметок -- на них можно будет вести учет своего золота, здоровья, состояний и предметов в инвентаре.

Стоит отметить, что правая колонка -- краткое описание персонажа, на которое я бы опирался, рассказывая его полную историю.

% }}}

% {{{ equipment


\end{multicols}
\chapter{Equipment}

\noindent
\rowcolors{1}{}{lightgray}
\setlength\tabcolsep{0pt}
\setlength\LTleft{0pt}
\setlength\LTright{0pt}
\begin{longtable}{p{0.4\textwidth}P{0.1\textwidth}p{0.4\textwidth}P{0.1\textwidth}}
  \textbf{\ \ \ name}
  & \textbf{price}
  & \textbf{\ \ \ name}
  & \textbf{price}
  \\
  \hline
  \ \ \ зелье здоровья                  &  25g &
  \ \ \ зелье языка животных            &  10g \\
  \ \ \ зелье ночного зрения            &  10g &
  \ \ \ зелье подводного дыхания        &  10g \\
  \ \ \ зелье паучьих ног               &  25g &
  \ \ \ фляска кислоты                  &  10g \\
  \ \ \ фляска вечной мерзлоты          &  10g &
  \ \ \ фляска дымовой завесы           &   5g \\
  \ \ \ фляска света                    &  10g &
  \ \ \ порошок усыпления               &   5g \\
  \ \ \ горсть пороха                   &   5g &
  \ \ \ палатка                         &  10g \\
  \ \ \ котел                           &   5g &
  \ \ \ огниво                          &   5g \\
  \ \ \ соль и перец                    &   2g &
  \ \ \ бинты                           &  15g \\
  \ \ \ бутылка виски                   &   3g &
  \ \ \ отмычки                         &   5g \\
  \ \ \ крюк                            &   5g &
  \ \ \ железные колючки                &   5g \\
  \ \ \ зеркало                         &  10g &
  \ \ \ факел                           &   3g \\
  \ \ \ веревка                         &   5g &
  \ \ \ мыло                            &   1g \\
  \ \ \ лопата                          &   5g &
  \ \ \ топорик                         &   5g \\
  \ \ \ кирка                           &   5g &
  \ \ \ музыкальный инструмент          &   5g
  \\

\end{longtable}

По поводу предметов -- спрашивайте на месте.

\begin{multicols}{2}

% }}}

% {{{ rules

\end{multicols}
\chapter{Rules}
\begin{multicols}{2}

\section{Ability Checks}

\lettrine{В}{}случае, когда действие игрока имеет шанс на провал, DM просит его сделать ролл на проверку аттрибута.

Успешность определяется сложением ролла d12 и мода, указанного DM.

Для успешного выполнения действия, итоговый ролл должен быть равен или превосходить сложность установленную DM. Сложность при этом не разглашается.

\section{Free Movement}

Вне боя игроки передвигаются свободно и взаимодействуют с окружением без какого либо установленного порядка.

\section{Combat}

Бой происходит по раундами. За один раунд каждый участник боя включая игроков и других персонажей ходит один раз.

\subsection{Initiative}

Если игроки устраивают засаду, у них есть возможность установить порядок ходов между собой.

В противном случае, порядок определяется роллами \textbf{d12 + dex mod} от большего к меньшему.

\subsection{Vision}

Персонажи могу взаимодействовать только с теми объектами, которые видят или чувствуют другим образом. Стены, дым, темнота и т.д. препятствуют зрению.

\subsection{Rolling}

Успешность каждого действия в бою определяется одним роллом d12. 

Урон, наносимый атакой или заклинанием равен сумме данного ролла и соответствующего мода. Исключения составляют критические провалы, когда все идет не по плану, и критические успехи. В случае с последними, урон удваивается.

\subsection{Opportunity Attacks}

Если персонаж, находясь в ближем бою с другим персонажем делает попытку сбежать, по нему проходит атака возможности. У данной ситуации есть исключения: критический успех побега, использование умения.

% 5 damage

\section{Resting}

На протяжение кампании у игроков будет возможность разбить лагерь и отдохнуть. Качество отдыха определяется действиями персонажей. Отдых восстанавливает игрокам здоровье.

Для хорошего отдыха необходимо хорошо поесть, обустроить лагерь и рассказать пару историй у костра.

% }}}

% {{{ print only

\end{multicols}
\pagebreak


\rowcolors{1}{}{}

\noindent
\setlength\tabcolsep{0pt}
\def\arraystretch{1.3}
\begin{tabularx}{\linewidth}{@{}YYYYYYYYYYYYY@{}}

  \multicolumn{13}{c}{\large\sc\textbf{Name:}} \\ 
  \multicolumn{13}{c}{\Large\emph{}} \\ \cline{4-10}

  \addlinespace[0.5cm]

  \multicolumn{6}{c}{\large\sc\textbf{Attributes:}} &
  \multicolumn{1}{c}{} &
  \multicolumn{6}{c}{\large\sc\textbf{Character:}}
  \\ \cline{1-6}

  \multicolumn{2}{c|}{\small\sc\textbf{Str}} &
  \multicolumn{2}{c|}{\small\sc\textbf{Dex}} &
  \multicolumn{2}{c}{\small\sc\textbf{Int}} &
  \multicolumn{1}{c}{} &
  \multicolumn{6}{Z}{\emph{}}
  \\ \cline{8-13}

  \multicolumn{2}{c|}{\LARGE{}} &
  \multicolumn{2}{c|}{\LARGE{}} &
  \multicolumn{2}{c}{\LARGE{}} &
  \multicolumn{1}{c}{} &
  \multicolumn{6}{Z}{\emph{}}
  \\ \cline{1-6} \cline{8-13}

  \multicolumn{2}{c|}{\large{}} & 
  \multicolumn{2}{c|}{\large{}} &
  \multicolumn{2}{c}{\large{}} &
  \multicolumn{1}{c}{} &
  \multicolumn{6}{Z}{\emph{}}
  \\ \cline{8-13}

  \multicolumn{6}{c}{\large\sc\textbf{Stats:}}  &
  \multicolumn{1}{c}{} &
  \multicolumn{6}{Z}{\emph{}}
  \\ \cline{1-6} \cline{8-13}

  \multicolumn{3}{c|}{\small\sc\textbf{Max Hp}} &
  \multicolumn{3}{c}{\small\sc\textbf{Gold}} &
  \multicolumn{1}{c}{} &
  \multicolumn{6}{Z}{\emph{}}
  \\ \cline{8-13}

  \multicolumn{3}{c|}{\LARGE{}} &
  \multicolumn{3}{c}{\LARGE{}} &
  \multicolumn{1}{c}{} &
  \multicolumn{6}{Z}{\emph{}}
  \\ \cline{1-6} \cline{8-13}
  &  &  &  &  &  &  &  &  &  &  &  & \\
\end{tabularx}

\vspace{40pt}

\noindent
\setlength\tabcolsep{0pt}
\def\arraystretch{1.3}
\begin{tabularx}{\linewidth}{@{}YYYYYYYYYYYYY@{}}

  \multicolumn{13}{c}{\large\sc\textbf{Name:}} \\ 
  \multicolumn{13}{c}{\Large\emph{}} \\ \cline{4-10}

  \addlinespace[0.5cm]

  \multicolumn{6}{c}{\large\sc\textbf{Attributes:}} &
  \multicolumn{1}{c}{} &
  \multicolumn{6}{c}{\large\sc\textbf{Character:}}
  \\ \cline{1-6}

  \multicolumn{2}{c|}{\small\sc\textbf{Str}} &
  \multicolumn{2}{c|}{\small\sc\textbf{Dex}} &
  \multicolumn{2}{c}{\small\sc\textbf{Int}} &
  \multicolumn{1}{c}{} &
  \multicolumn{6}{Z}{\emph{}}
  \\ \cline{8-13}

  \multicolumn{2}{c|}{\LARGE{}} &
  \multicolumn{2}{c|}{\LARGE{}} &
  \multicolumn{2}{c}{\LARGE{}} &
  \multicolumn{1}{c}{} &
  \multicolumn{6}{Z}{\emph{}}
  \\ \cline{1-6} \cline{8-13}

  \multicolumn{2}{c|}{\large{}} & 
  \multicolumn{2}{c|}{\large{}} &
  \multicolumn{2}{c}{\large{}} &
  \multicolumn{1}{c}{} &
  \multicolumn{6}{Z}{\emph{}}
  \\ \cline{8-13}

  \multicolumn{6}{c}{\large\sc\textbf{Stats:}}  &
  \multicolumn{1}{c}{} &
  \multicolumn{6}{Z}{\emph{}}
  \\ \cline{1-6} \cline{8-13}

  \multicolumn{3}{c|}{\small\sc\textbf{Max Hp}} &
  \multicolumn{3}{c}{\small\sc\textbf{Gold}} &
  \multicolumn{1}{c}{} &
  \multicolumn{6}{Z}{\emph{}}
  \\ \cline{8-13}

  \multicolumn{3}{c|}{\LARGE{}} &
  \multicolumn{3}{c}{\LARGE{}} &
  \multicolumn{1}{c}{} &
  \multicolumn{6}{Z}{\emph{}}
  \\ \cline{1-6} \cline{8-13}
  &  &  &  &  &  &  &  &  &  &  &  & \\
\end{tabularx}

% }}}

\end{document}
